\section{Metodología}
\label{sec:metodologia}

El análisis compara dos versiones de cada enzima:

\begin{enumerate}
    \item WT: enzima original utilizada como referencia;
    \item ENG: enzima modificada mediante ingeniería.
\end{enumerate}

El objetivo es determinar si ENG dura más, convierte mejor y termina generando
más producto.

\subsection{Diseño del Análisis}
\label{subsec:diseno}

Se evaluaron cuatro aplicaciones industriales:

\begin{enumerate}
    \item proteasa para procesamiento de alimentos;
    \item beta-glucosidasa para aprovechamiento de bagazo de caña;
    \item xilanasa para procesamiento de biomasa;
    \item beta-glucosidasa para síntesis de glucósidos.
\end{enumerate}

Los casos representan restricciones industriales diferentes y cuentan con
antecedentes de ingeniería orientada a mejorar estabilidad, actividad o ambas
propiedades
\citep{yuan_apre2709_2025,cao_ks5a7_2018,zhang_xylanase_2010,yin_bgl1d_2019}.

Cada condición se evaluó con cuatro réplicas.

\subsection{Protocolo y Medidas}
\label{subsec:protocolo}

\begin{table}[H]
\centering
\tabcaptionlink{anx:tab_protocolo}{Protocolo Utilizado para Comparar WT y ENG}{tab:protocolo_metodologia}



\renewcommand{\arraystretch}{1.3}

\begin{tabularx}{0.96\textwidth}{
|p{2.6cm}|
p{4.1cm}|
X|
}
\hline
Prueba & Condición & Resultado Utilizado \\ \hline

Estabilidad &
\(60\,^{\circ}\mathrm{C}\), 10 mM de sustrato &
\(t_{50}\) y actividad conservada en el tiempo. \\ \hline

Cinética &
0.5--80 mM de sustrato &
Eficiencia de conversión e inhibición por exceso de sustrato. \\ \hline

Respuesta Térmica &
\(35--90\,^{\circ}\mathrm{C}\) &
Temperatura de máxima actividad. \\ \hline

Producción Base &
\(60\,^{\circ}\mathrm{C}\), 40 mM, 8 horas &
Producto acumulado bajo una condición común. \\ \hline

Exploración Operativa &
Misma malla de temperaturas y sustratos, 8 horas &
Mejor punto de producción encontrado para WT y ENG. \\ \hline

\end{tabularx}
\end{table}

La condición base permite comparar todas las variantes bajo el mismo entorno.
La exploración operativa se utiliza después para responder una pregunta
distinta: qué combinación de temperatura y sustrato conviene usar para cada
variante dentro del rango evaluado.

\subsection{Indicadores de Comparación}
\label{subsec:indicadores}

Se utilizan tres índices relativos:

\begin{enumerate}
    \item Índice de estabilidad (\(I_S\)): compara la actividad conservada por
    ENG frente a WT.

    \item Índice de eficiencia (\(I_C\)): compara la eficiencia catalítica de
    ENG frente a WT.

    \item Índice de producción (\(I_P\)): compara el producto acumulado por ENG
    frente a WT bajo la condición base.
\end{enumerate}

En los tres casos:

\begin{equation}
\text{Índice}
=
\frac{\text{Resultado de ENG}}
{\text{Resultado de WT}}.
\end{equation}

La interpretación es directa:

\begin{table}[H]
\centering
\tabcaptionlink{anx:tab_lectura_indices}{Lectura de los Índices}{tab:lectura_indices}



\renewcommand{\arraystretch}{1.25}

\begin{tabular}{|c|p{9.5cm}|}
\hline
Valor & Interpretación \\ \hline
\(>1\) & ENG mejora el resultado frente a WT. \\ \hline
\(=1\) & No existe diferencia entre ambas variantes. \\ \hline
\(<1\) & ENG obtiene un resultado inferior a WT. \\ \hline
\end{tabular}
\end{table}

También se utiliza el factor de traducción productiva (\(F_P\)):

\begin{equation}
F_P=
\frac{I_P}{I_S I_C}.
\end{equation}

Si \(F_P<1\), parte de la mejora molecular es absorbida por otras restricciones;
si \(F_P>1\), factores adicionales amplifican la producción.

\subsection{Restricciones Consideradas}
\label{subsec:restricciones_metodologia}

El análisis incorpora cuatro mecanismos que pueden cambiar la decisión:

\begin{enumerate}
    \item pérdida de actividad con el tiempo;
    \item inhibición por exceso de sustrato;
    \item inhibición por producto acumulado;
    \item cambios en selectividad hacia el producto deseado.
\end{enumerate}

Las funciones utilizadas para representar estos mecanismos se presentan en el
\hyperref[sec:reproducibilidad]{Anexo C}.

\subsection{Flujo de Decisión}
\label{subsec:flujo}

\begin{figure}[H]
\centering

\begin{tikzpicture}[
    node distance=0.72cm,
    caja/.style={
        draw,
        rounded corners=4pt,
        minimum width=9.2cm,
        minimum height=0.88cm,
        align=center
    },
    >=Latex
]

\node[caja] (comparacion)
{Comparar WT frente a ENG};

\node[caja, below=of comparacion] (medicion)
{Medir estabilidad, eficiencia, temperatura y producción};

\node[caja, below=of medicion] (indices)
{Cuantificar la mejora relativa};

\node[caja, below=of indices] (restriccion)
{Identificar el cuello de botella};

\node[caja, below=of restriccion] (decision)
{Elegir la condición que maximiza producto útil};

\draw[->, thick] (comparacion) -- (medicion);
\draw[->, thick] (medicion) -- (indices);
\draw[->, thick] (indices) -- (restriccion);
\draw[->, thick] (restriccion) -- (decision);

\end{tikzpicture}

\figcaptionlink{anx:fig_flujo}{Flujo Metodológico Orientado a Decisión}{fig:flujo_metodologia}


\end{figure}

La \hyperref[sec:evidencia]{Sección 4} presenta los resultados bajo condición
común y la \hyperref[sec:discusion]{Sección 5} utiliza esos resultados para
definir la estrategia de explotación de cada aplicación.
